\documentclass[11pt]{article}
\usepackage[utf8]{inputenc}
\usepackage[T1]{fontenc}
\usepackage[french]{babel}
\usepackage{graphicx}
\usepackage{booktabs}
\usepackage{geometry}
\geometry{margin=2.2cm}

\title{Détection de password spraying par extraction de features et Isolation Forest}
\author{\textit{(Ton nom)} \\
ENSIMAG}
\date{\today}

\begin{document}
\maketitle

\section{Objectif}
Ce projet vise à détecter des anomalies d'authentification compatibles avec des attaques de type \emph{password spraying} à partir de logs.

\section{Données}
Décrire: source des logs, période, champs disponibles, pré-traitement.

\section{Features}
Décrire: n\_attempts, n\_fail, n\_users, fail\_rate, attempts\_per\_min, entropy, etc.

\section{Modèle}
Décrire: IsolationForest (non supervisé), contamination, standardisation.

\section{Résultats}
\subsection{Top alertes}
\input{../data/processed/tables/alerts_top.tex}

\subsection{Figures}
\begin{figure}[h]
\centering
\includegraphics[width=0.85\linewidth]{../data/processed/figures/anomaly_score_hist.png}
\caption{Distribution des scores d'anomalie}
\end{figure}

\begin{figure}[h]
\centering
\includegraphics[width=0.95\linewidth]{../data/processed/figures/anomalies_over_time.png}
\caption{Anomalies détectées dans le temps}
\end{figure}

\section{Limites et améliorations}
Expliquer: fenêtres non-overlapping vs sliding, distinct counts, interprétabilité, évaluation avec labels.

\end{document}
